\section{Removing the Critical Section}\label{sec:batching}

The analysis at the end of \cref{sec:approach} predicts where the fused engine fails, and
profiling confirms it. On Fischer $N{=}10$ with twelve threads, a sampling profile attributes
$6\,049$ samples to the insertion routine itself and $4\,949$ to zone interning, but roughly
$16\,000$ to the four mutex symbols --- acquiring, waiting on, and releasing stripe locks.
More time is spent \emph{waiting} for antichains than traversing them. Fischer concentrates
its symbolic states on few configurations, so all twelve threads converge on a handful of
stripes, and each holds its lock for a traversal that grows as the search proceeds. This is
the same dependence on $\rho$ that prior work reports; the difference is what we do about it.

The prior remedy weakens the data structure, replacing deletion-then-insertion with an atomic
replacement so that threads no longer block, at the price of losing the invariant that the
stored set is an antichain and thus allowing $|\Sigma|$ to grow. We keep the structure exactly
as it is and change \emph{when} the work happens.

Observe that nothing in \cref{alg:insert} needs to happen at the moment a successor is
generated. Its outcome depends only on the antichain and the incoming zone, so it can be
deferred to the end of the level, provided every successor is remembered in the meantime. Our
second engine does precisely that. During expansion a thread appends the successor to its
configuration's chain and nothing more: acquire, push, release. No comparison happens under
the lock, so the critical section is $O(1)$ irrespective of $\rho$. Once expansion is
complete, every configuration touched during the level is \emph{settled}: its chain is
partitioned into the antichain accumulated by earlier levels and the zones appended during
this one, and the pending zones are folded in by exactly the steps of \cref{alg:insert} ---
with no lock held at all.

Settling is lock-free by construction rather than by convention. A settling thread modifies
only the \texttt{dead} and zone-identifier fields of nodes on the chain it owns, and never
appends to the arena or grows the table, so two configurations that happen to share a stripe
cannot interfere. \Cref{fig:rho} shows why contention tracks $\rho$ at all, and \cref{fig:timeline} the resulting thread schedules.

\begin{figure}[t]
\centering
\begin{tikzpicture}[font=\small, x=1mm, y=1mm,
  wrk/.style={draw=black!70, fill=black!10, minimum width=8mm, minimum height=6mm, inner sep=0pt},
  zn/.style={draw=black!60, fill=white, minimum width=7mm, minimum height=3.2mm, inner sep=0pt},
  ttl/.style={font=\bfseries\small},
  sub/.style={font=\scriptsize},
  a/.style={-{Latex[length=1.5mm]}, draw=black!70},
  bar/.style={draw=black!70, fill=black!35, minimum height=3mm, inner sep=0pt}]

% ================= LEFT: many thin antichains =================
\node[ttl] at (24,8) {many thin antichains};
\node[sub] at (24,2) {$\rho \approx 1$ \;(Viking)};
\foreach \i/\x in {0/4, 1/17, 2/30, 3/43} \node[wrk] at (\x,-8) {$T_\i$};
\foreach \i/\x in {0/4, 1/17, 2/30, 3/43} {
  \draw[a] (\x,-11.5) -- (\x,-17);
  \node[zn] at (\x,-20) {};
}
\foreach \x in {10, 23, 36} \node[zn, draw=black!35] at (\x,-20) {};
\node[sub, anchor=east] at (-2,-20) {visited};
\node[sub] at (24,-27) {every thread lands on its own antichain};
\node[bar, minimum width=5mm, anchor=west] at (4,-35) {};
\node[bar, minimum width=5mm, anchor=west] at (17,-35) {};
\node[bar, minimum width=5mm, anchor=west] at (30,-35) {};
\node[bar, minimum width=5mm, anchor=west] at (43,-35) {};
\node[sub, anchor=east] at (-2,-35) {lock held};

% ================= RIGHT: few fat antichains =================
\node[ttl] at (95,8) {few fat antichains};
\node[sub] at (95,2) {$\rho \gg 1$ \;(Fischer)};
\foreach \i/\x in {0/75, 1/88, 2/101, 3/114} \node[wrk] at (\x,-8) {$T_\i$};
\foreach \x in {75, 88, 101, 114} \draw[a] (\x,-11.5) -- (95,-16.5);
\foreach \k in {0,...,5} \node[zn] at (95,-19.5-3.6*\k) {};
\node[sub] at (95,-45) {all four contend for one antichain};
\node[bar, minimum width=26mm, anchor=west] at (75,-52) {};
\node[sub, anchor=east] at (73,-52) {lock held};
\end{tikzpicture}
\caption{Why $\rho$ decides contention. Where a configuration carries a single zone, threads
touch distinct antichains and each holds its stripe lock for one comparison. Where a
configuration accumulates many zones, threads converge on the same antichain and each holds the
lock across a traversal proportional to its length, so the traversals are executed serially.
The batched engine performs exactly the same traversals, but after the lock is released.}
\label{fig:rho}
\end{figure}

\begin{figure*}[t]
\centering
\begin{tikzpicture}[font=\small, x=1mm, y=1mm,
  wait/.style={draw=black!50, fill=black!10, minimum height=5mm, inner sep=0pt},
  work/.style={draw=black!75, fill=black!45, minimum height=5mm, inner sep=0pt, text=white},
  free/.style={draw=black!60, pattern=north east lines, pattern color=black!45,
               minimum height=5mm, inner sep=0pt}]
\node[lbl,anchor=east] at (-3,14) {\textbf{(a) fused}};
\foreach \t/\off in {1/0, 2/13, 3/26} {
  \node[lbl,anchor=east] at (-3,14-6*\t) {$T_\t$};
  \node[wait,minimum width=\off mm,anchor=west] at (0,14-6*\t) {};
  \node[work,minimum width=13mm,anchor=west] at (\off,14-6*\t) {\scriptsize scan};
  \node[wait,minimum width=39mm,anchor=west] at (\off+13,14-6*\t) {};
  \node[work,minimum width=13mm,anchor=west] at (\off+52,14-6*\t) {\scriptsize scan};
}
\node[lbl,anchor=west] at (82,2) {\scriptsize the $\Theta(\rho)$ scans execute \emph{serially}};
\node[lbl,anchor=east] at (-3,-10) {\textbf{(b) batched}};
\foreach \t/\off in {1/0, 2/2, 3/4} {
  \node[lbl,anchor=east] at (-3,-10-6*\t) {$T_\t$};
  \node[work,minimum width=2mm,anchor=west] at (\off,-10-6*\t) {};
  \node[work,minimum width=2mm,anchor=west] at (\off+4,-10-6*\t) {};
  \node[free,minimum width=52mm,anchor=west] at (10,-10-6*\t) {\scriptsize settle};
}
\draw[dashed,black!55] (8,-13) -- (8,-30);
\node[lbl,anchor=north] at (8,-31) {\scriptsize expansion ends};
\node[lbl,anchor=west] at (82,-22) {\scriptsize the same scans execute \emph{concurrently}};
\node[work,minimum width=6mm] at (10,26) {}; \node[lbl,anchor=west] at (14,26) {\scriptsize holds the stripe lock};
\node[wait,minimum width=6mm] at (52,26) {}; \node[lbl,anchor=west] at (56,26) {\scriptsize blocked on it};
\node[free,minimum width=6mm] at (92,26) {}; \node[lbl,anchor=west] at (96,26) {\scriptsize antichain fold, no lock};
\end{tikzpicture}
\caption{Thread schedules on a configuration with a long antichain. Under fusion (a) each
insertion holds the stripe lock across a $\Theta(\rho)$ traversal, so threads queue and the
traversals are effectively serialised. Under batching (b) insertions are constant-time pushes
that barely contend, and the traversals --- now one per configuration per level rather than
one per successor --- run concurrently outside any lock.}
\label{fig:timeline}
\end{figure*}

Two implementation details are load-bearing. The first is that the fold must run over a
compacted array of live nodes rather than over the chain itself. Deferring subsumption lets a
configuration accumulate thousands of pending nodes, and a node flagged dead is skipped but
still \emph{traversed}, so folding over the chain is quadratic in the number of pending zones;
measured on Fischer $N{=}10$, the chain-walking version took $8.17$\,s against $2.92$\,s for
the fused engine it was meant to improve, while the compacted version is faster than both. The
second is that the $\aLU{}$ bounds of line~\ref{ln:bounds}, which depend only on the
configuration, are now computed once per configuration per level instead of once per
successor --- an optimisation the fused engine cannot make, because it never sees a
configuration's successors together.

Identifying which configurations were touched costs nothing. Because \cref{alg:insert}
prepends, the head of a chain is always its most recent node, so its level stamp answers
``has this configuration already been touched during this level?'' under the lock we are
already holding. Each touched configuration is thus recorded exactly once, with no auxiliary
set and no additional synchronisation.

The two engines are complementary, and $\rho$ decides between them. Where $\rho \approx 1$
antichains are singletons, the fused critical section is already minimal, and batching adds
bookkeeping for nothing. Where $\rho$ is large the fused engine serialises on a few stripes
while the batched engine performs identical comparisons in parallel. Crucially this is a
property of the model rather than of the tool, so it can be estimated in advance --- a few
hundred levels of exploration suffice --- and used to choose an engine before committing to a
long run.
